\documentclass[11pt]{article}
\usepackage[spanish]{babel}
\usepackage[margin=1in]{geometry}
\usepackage{amsmath, amssymb, amsfonts}
\usepackage{booktabs}
\usepackage{enumitem}
\usepackage{tcolorbox}
\usepackage{hyperref}

\title{\textbf{Hoja de Trabajo 1: Agentes y Entornos} \vspace{1em}\\
Inteligencia Artificial - CC3085}
\author{José Antonio Mérida Castejón}
\date{\today}

\begin{document}
\maketitle

\section*{Task 1: Teoría}
\noindent\textit{Responda con criterio y análisis de ingenierx. No se esperan definiciones de libro, sino que analice las
consecuencias de las decisiones de diseño.}

\subsection*{La Falacia del Jugador y la Propiedad de Markov}
\noindent\textit{Durante la clase definimos que un MDP cumple con la propiedad de Markov. Con esto en mente
 considere que usted está diseñando un agente para jugar Blackjack. Un ``experto'' le dice que debe
 incluir en el estados el historial de las últimas 10 manos ganadas/perdidas porque ``si ha perdido
 10 veces seguidas, le toca ganar''. Con esto responda}

 \begin{itemize}
 \item\textit{Explique matemáticamente por qué incluir el historial de victorias pasadas viola (o es irrelevante
para) la propiedad de Markov si el juego es verdaderamente aleatorio. ¿Qué información sí es
crítica incluir en el estado s para que el MDP sea válido?}
\end{itemize}

La propiedad de Markov funciona dado que los eventos son independientes, dónde el resultado de tomar
una decisión o llevar a cabo una acción depende únicamente del estado actual en el que estamos. Es decir,
el pasado pudo habernos llevado al estado actual, pero sin importar como nosotros hayamos llegado 
únicamente importa la situación actual.

\subsection*{El Dilema del Factor Descuento ($\gamma$)}
\noindent\textit{Considere lo visto en clase y además tome en cuenta dos escenarios.}

\begin{itemize}
  \item\textbf{Escenario A:} Un robot de rescate en un edificio en llamas.
  \item\textbf{Escenario B:} Un agente de inversión para un fondo de pensiones a 30 años.
\end{itemize}

\textit{Con ello responda:}

\begin{itemize}
  \item \textit{¿Qué valor de γ (cercano a 0 o cercano a 1) asignaría a cada escenario? Justifique su respuesta
explicando cómo el valor de γ afecta la "impaciencia" o "visión a largo plazo" del agente y las
consecuencias de elegir el valor incorrecto para cada caso.}
\end{itemize}

Cada uno de estos ejemplos presenta un extremo en cuanto a prioridad de eventos \textit{cercanos} vs
\textit{a futuro}. En en \textbf{Escenario A}, el tiempo es clave y el priorizar alcanzar a una alta
cantidad de gente puede resultar en evitar el rescate de personas que se encuentren en una 
situación menos peligrosa. Si priorizamos beneficios tardíos, estaríamos entrenando al robot para
salvar personas que ya no pueden ser salvadas. Por otro lado, en el \textbf{Escenario B} debemos
utilizar el \textit{approach} contrario. Aquí deberíamos de priorizar los retornos a largo plazo, ya 
que de momento no nos es útil en lo absoluto beneficiarnos hasta que el fondo de pensiones se
ponga en uso.

\subsection*{Políticas Estocásticas vs Deterministas}
\textit{Considere que durante la clase dijimos que una solución es una política $\pi(s)$. Con eso en mente responda:}

\begin{enumerate}
  \item\textit{Si el entorno es estocástico (mis acciones fallan a veces), ¿necesito una política estocástica (que a
    veces elija A y a veces B para el mismo estado) para ser óptimo, o basta con una política
  determinista? Justifique su respuesta basándose en la ecuación de Bellman para $V_{opt}$}

  No es necesario emplear una política estocástica, en este caso se evalúa el retorno \textit{esperado}
  de tomar una acción determinada en un estado determinado. Por ejemplo, si tuviésemos una elección
  de tirar una moneda y ganar $Q50,000$ en cara y perder $Q5$ en cruz o simplemente no tirarla. No
  sería lógico para cualquier persona razonable evitar tirar la moneda, \textit{``Como hay veces
  que pierdo $Q5$, esta vez no la voy a tirar''} - nadie nunca.
\end{enumerate}

\end{document}
